\chapter{Espaços de Hilbert}\label{ch:b3:hilbert}

Um \omterm{def:b3:hilbert:inner}{espaço de Hilbert} é um espaço de Banach cuja norma provém de um
\omterm{def:b3:hilbert:inner}{produto interno} --- e essa única estrutura adicional restaura, em
dimensão infinita, quase toda a geometria euclidiana: existem projeções
ortogonais, todo funcional \omterm{def:b3:topology:continuity}{contínuo} é um \omterm{def:b3:hilbert:inner}{produto interno} contra um
vetor fixado (Riesz), e as bases ortonormais expandem todo vetor em uma
série convergente com contabilidade pitagórica (\omterm{thm:b3:hilbert:parseval}{Parseval}). O ponto alto
do capítulo é uma dívida honrada: o sistema trigonométrico é uma base
ortonormal de $L^2$, de modo que a identidade de \omterm{thm:b3:hilbert:parseval}{Parseval} vale para
\emph{toda} função de quadrado \omterm{def:b3:lebesgue:l1}{integrável} --- o enunciado que o segundo
ano só pôde demonstrar para funções $\mathcal C^1$ por partes. Encerramos com
Lax--Milgram, o lema cavalo de batalha da abordagem variacional das
equações diferenciais.

Ao longo de todo o capítulo, $H$ é um espaço vetorial sobre
$K = \R$ ou $\C$.

\section{Produtos internos; o teorema da projeção}

\begin{definition}\label{def:b3:hilbert:inner}
Um \emph{produto interno}\index{produto interno} é uma aplicação $\langle
\cdot,\cdot\rangle \colon H\times H \to K$,
linear na segunda variável, com $\langle y, x\rangle =
\overline{\langle x, y\rangle}$ e $\langle x, x\rangle > 0$ para $x \neq 0$. Ele induz a norma
$\norm x = \langle x,
x\rangle^{1/2}$, a \emph{desigualdade de Cauchy--Schwarz} $\abs{\langle x, y\rangle} \leq \norm x\norm y$ (a demonstração do
segundo ano --- o discriminante --- não muda) e a \emph{lei do
paralelogramo}
\[
\norm{x + y}^2 + \norm{x - y}^2 = 2\norm x^2 + 2\norm y^2 .
\]
Um \emph{espaço de Hilbert}\index{espaço de Hilbert} é um espaço com
produto interno \omterm{def:b3:complete:complete}{completo} para essa norma. Exemplos: $\ell^2$
(\cref{pb:b3:banach:1}) e, o fundamental, $L^2(\mu)$ com $\langle f, g\rangle = \int\bar fg\,\dd\mu$ --- \omterm{def:b3:complete:complete}{completo} por
Riesz--Fischer (\cref{thm:b3:lp:complete}); o produto interno é finito por
Cauchy--Schwarz ($= $ Hölder em $p = q = 2$).
\end{definition}

\begin{theorem}[Projeção sobre um convexo fechado]
\label{thm:b3:hilbert:projection}
Sejam $C \neq \varnothing$ um subconjunto \emph{convexo} fechado do \omterm{def:b3:hilbert:inner}{espaço de Hilbert}
$H$ e $x \in H$. Existe um único $p_C(x)
\in C$ com\index{teorema da projeção}
\[
\norm{x - p_C(x)} = d(x, C),
\]
caracterizado por: $\operatorname{Re}\langle x - p_C(x),\ c -
p_C(x)\rangle \leq 0$ para todo $c \in C$. A aplicação $p_C$ é
$1$-lipschitziana.
\end{theorem}

\begin{proof}
Sejam $d = d(x, C)$ e $(c_n) \subseteq C$ com $\norm{x - c_n}
\to d$. A lei do paralelogramo em $x - c_n$ e $x - c_m$:
\[
\norm{c_n - c_m}^2 = 2\norm{x - c_n}^2 + 2\norm{x - c_m}^2 -
4\,\bigl\|x - \tfrac{c_n + c_m}2\bigr\|^2
\leq 2\norm{x{-}c_n}^2 + 2\norm{x{-}c_m}^2 - 4d^2
\]
(a convexidade coloca o ponto médio em $C$): o membro direito tende a
$0$, de modo que $(c_n)$ é de Cauchy, e seu limite $p \in C$ (fechado)
atinge $d$. Unicidade: dois minimizadores dão, pela mesma
identidade, $\norm{p - p'}^2 \leq 2d^2 + 2d^2 - 4d^2 = 0$.

Caracterização: para $c \in C$, $t \in \intoc01$, o vetor $p + t(c - p) \in C$, de modo que
\[
d^2 \leq \norm{x - p - t(c-p)}^2
= d^2 - 2t\operatorname{Re}\langle x - p, c - p\rangle +
t^2\norm{c-p}^2 ;
\]
divida por $t \to 0^+$: $\operatorname{Re}\langle x - p, c -
p\rangle \leq 0$. Reciprocamente, essa desigualdade dá $\norm{x -
c}^2 = \norm{x - p}^2 - 2\operatorname{Re}\langle x - p, c -
p\rangle + \norm{p - c}^2 \geq \norm{x-p}^2$.
Lipschitz: para $x, y$ de projeções $p, q$, some as duas
desigualdades variacionais (com $c = q$ e com $c = p$,
respectivamente): $\operatorname{Re}\langle x - y - (p - q), p - q\rangle \geq
0$, logo $\norm{p - q}^2 \leq \operatorname{Re}\langle x - y, p -
q\rangle \leq \norm{x - y}\norm{p - q}$.
\end{proof}

\begin{theorem}[Decomposição ortogonal]
\label{thm:b3:hilbert:decomposition}
Seja $F$ um \emph{subespaço fechado} de $H$. Então $p_F$ é
linear, $x - p_F(x) \perp F$ para todo $x$, e\index{complemento ortogonal}
\[
H = F \oplus F^\perp,
\qquad F^\perp = \{y : \langle y, f\rangle = 0\ \forall f\in
F\},
\qquad (F^\perp)^\perp = F .
\]
Para um subespaço geral, $(F^\perp)^\perp = \bar F$; em particular, $F$ é denso se, e
somente se, $F^\perp = \{0\}$.
\end{theorem}

\begin{proof}
Para um subespaço, a caracterização variacional com $c =
p_F(x) \pm f$ ($f \in F$, nos dois
sinais, e $\iu f$ no caso complexo) força $\langle x - p_F(x), f\rangle = 0$: o resíduo é ortogonal
a $F$. Decomposição $x = p_F(x) + (x
- p_F(x))$ com $F \cap F^\perp = \{0\}$ ($\langle y, y
\rangle = 0$); a linearidade de $p_F$
decorre da unicidade de tais decomposições (ambos os lados são lineares
nelas). $(F^\perp)
^\perp \supseteq F$ sempre; reciprocamente, se $x \perp F^\perp$, escreva $x = f + g$: $g = x - f \in F^\perp$ e
$\langle g,
g\rangle = \langle x, g\rangle - \langle f, g\rangle = 0$: $x
= f \in F$. Para um subespaço geral $F$: $F^\perp = \bar
F^{\,\perp}$ (\omterm{def:b3:topology:continuity}{continuidade} do
\omterm{def:b3:hilbert:inner}{produto interno}), de modo que $(F^\perp)^\perp = \bar F$ pelo caso fechado; e a \omterm{ex:b3:lebesgue:gamma}{densidade}
equivale a $\bar F = H$, que equivale a $F^\perp = 0$.
\end{proof}

\begin{theorem}[Representação de Riesz]\label{thm:b3:hilbert:riesz}
Para todo funcional linear \omterm{def:b3:topology:continuity}{contínuo} $\varphi \in H'$ existe um único $a \in H$
com\index{teorema de representação de Riesz}
\[
\varphi(x) = \langle a, x\rangle \quad (x \in H),
\qquad \norm\varphi_{H'} = \norm a .
\]
\end{theorem}

\begin{proof}
Se $\varphi = 0$: $a = 0$. Do contrário, $F = \ker\varphi$ é um subespaço próprio fechado;
tome $u \in F^\perp$, $\norm u = 1$ (\cref{thm:b3:hilbert:decomposition}: $F^\perp \neq 0$, pois $F \neq H$). Para todo $x$, o
vetor $\varphi(x)u -
\varphi(u)x \in \ker\varphi$, logo $\perp u$:
\[
0 = \langle u, \varphi(x)u - \varphi(u)x\rangle
= \varphi(x) - \varphi(u)\langle u, x\rangle :
\qquad \varphi(x) = \langle
\overline{\varphi(u)}\,u,\ x\rangle .
\]
Assim $a = \overline{\varphi(u)}u$ resolve. Unicidade: $\langle a
- a', x\rangle = 0$ para todo $x$; teste $x = a - a'$.
Normas: $\abs{\varphi(x)} \leq \norm a\norm x$ (Cauchy--Schwarz), com igualdade em $x = a$.
\end{proof}

\begin{example}[Uma projeção, calculada até o fim]
\label{ex:b3:hilbert:projectionexample}
Em $H = L^2(\intcc01)$, qual é a melhor aproximação de $f(x) = x^2$ por uma função afim? O
subespaço $F =
\operatorname{Vect}(1, x)$ é fechado (tem dimensão finita), e $p_F(f) = a + bx$ fica
caracterizada pela ortogonalidade do resíduo a $1$ e a $x$:
\[
\int_0^1(x^2 - a - bx)\,\dd x = 0,
\qquad
\int_0^1x\,(x^2 - a - bx)\,\dd x = 0,
\]
isto é, $\frac13 = a + \frac b2$ e $\frac14 = \frac a2 +
\frac b3$: $a = -\frac16$, $b = 1$. Logo $p_F(x^2) = x -
\frac16$, e o erro é
\[
d(f, F)^2 = \int_0^1\Bigl(x^2 - x + \frac16\Bigr)^2\dd x =
\frac1{180},
\qquad d(f, F) = \frac1{6\sqrt5} .
\]
Duas observações que vale internalizar. Primeira: o cálculo não passa
de um sistema linear $2\times2$ --- as \emph{equações normais}; para a base
monomial, sua matriz $\bigl(\frac1{i+j+1}\bigr)$ é a notoriamente mal condicionada matriz de
Hilbert, e ortogonalizar antes (polinômios de Legendre,
\cref{pb:b3:hilbert:1}) é o remédio. Segunda: a melhor aproximação
\emph{uniforme} de $x^2$ por funções afins é outra ($x - \frac18$, por
equioscilação): cada norma tem sua própria geometria, e só a
hilbertiana responde com um sistema linear.
\end{example}

\section{Bases ortonormais}

\begin{definition}\label{def:b3:hilbert:onb}
Uma família $(e_i)_{i\in I}$ é \emph{ortonormal} se $\langle
e_i, e_j\rangle = \delta_{ij}$, e é uma \emph{base
hilbertiana}\index{base hilbertiana} (base ortonormal) se, além disso,
suas combinações lineares finitas são densas em $H$ (a família é
\emph{total}). Tratamos o caso enumerável $I = \N$, que, por
Gram--Schmidt, cobre todo $H$ \emph{separável}
(\cref{prop:b3:hilbert:gramschmidt}).
\end{definition}

\begin{theorem}[Bessel, Parseval]\label{thm:b3:hilbert:parseval}
Sejam $(e_n)_{n\in\N}$ ortonormal em $H$ e $c_n(x) =
\langle e_n, x\rangle$.
\begin{enumerate}
  \item (Bessel) $\sum_n\abs{c_n(x)}^2 \leq \norm x^2$, e a série $\sum_nc_n(x)e_n$ converge em $H$, com soma
        $p_F(x)$, $F = \overline{\operatorname{Vect}}(e_n)$.
  \item São equivalentes: (i) $(e_n)$ é uma \omterm{def:b3:hilbert:onb}{base hilbertiana}; (ii)
        $x = \sum_nc_n(x)e_n$ para todo $x$; (iii) \emph{Parseval}\index{identidade
        de Parseval}: $\norm x^2
        = \sum_n\abs{c_n(x)}^2$ para todo $x$; (iv) o único vetor
        ortogonal a todos os $e_n$ é $0$.
  \item Se $(e_n)$ é uma \omterm{def:b3:hilbert:onb}{base hilbertiana}, $x \mapsto (c_n(x))_n$ é um isomorfismo
        isométrico $H \to \ell^2$ (\emph{todo} \omterm{def:b3:hilbert:inner}{espaço de Hilbert} separável de
        dimensão infinita ``é'' $\ell^2$), e $\langle x, y\rangle =
        \sum_n\overline{c_n(x)}c_n(y)$.
\end{enumerate}
\end{theorem}

\begin{proof}
(1) Para $N$ finito: $x - \sum_{n\leq N}c_ne_n \perp e_k$ ($k
\leq N$), de modo que Pitágoras dá $\norm x^2 = \sum_{n\leq
N}\abs{c_n}^2 + \norm{x - \sum_{n\leq N}c_ne_n}^2$:
eis Bessel. As somas parciais $S_N = \sum_{n\leq N}c_ne_n$ são de Cauchy: $\norm{S_N - S_M}^2 = \sum_{M<n\leq N}\abs{c_n}^2$, cauda de uma
série convergente; o limite está em $F$, e $x - \lim S_N
\perp$ cada $e_k$
(\omterm{def:b3:topology:continuity}{continuidade}), logo $\perp F$: pela unicidade da decomposição ortogonal,
$\lim S_N = p_F(x)$.

(2) (i)$\Rightarrow$(ii): $F = H$, logo $p_F = \mathrm{id}$. (ii)$\Rightarrow$(iii): Pitágoras no
limite ($\norm{S_N}^2
= \sum_{n \leq N}\abs{c_n}^2 \to \norm x^2$). (iii)$\Rightarrow$(iv): $x \perp$ a todos os $e_n$ dá $\norm x^2 =
0$.
(iv)$\Rightarrow$(i): $F^\perp = \{0\}$ (ser ortogonal a todos os $e_n$ é ser ortogonal a
$F$), de modo que $F$ é denso pelo~\cref{thm:b3:hilbert:decomposition}; mas $F$,
sendo um \omterm{def:b3:topology:interior}{fecho}, já é fechado: $F = H$.

(3) A aplicação é linear, isométrica por (iii) (logo injetora) e
sobrejetora: dado $(c_n) \in \ell^2$, a série $\sum
c_ne_n$ converge (de Cauchy como em (1))
para uma pré-imagem. A fórmula do \omterm{def:b3:hilbert:inner}{produto interno} vem por polarização a
partir de (iii), ou de um cálculo direto de limite.
\end{proof}

\begin{proposition}[Gram--Schmidt]\label{prop:b3:hilbert:gramschmidt}
Seja $(x_n)$ uma sequência linearmente independente. Pondo
indutivamente $\tilde e_n = x_n - \sum_{k<n}\langle e_k,
x_n\rangle e_k$ e $e_n = \tilde e_n/\norm{\tilde e_n}$, obtém-se uma família ortonormal $(e_n)$
com os mesmos subespaços gerados finitos: $\operatorname{Vect}(e_1, \dots, e_n) = \operatorname{Vect}
(x_1, \dots, x_n)$. Consequentemente, todo
\omterm{def:b3:hilbert:inner}{espaço de Hilbert} separável (isto é, com um subconjunto denso
enumerável) tem \omterm{def:b3:hilbert:onb}{base hilbertiana}.\index{ortonormalização de
Gram--Schmidt}
\end{proposition}

\begin{proof}
Indução: $\tilde e_n \perp e_k$ ($k < n$) por construção, e $\tilde e_n \neq 0$ pela independência;
os subespaços gerados coincidem em cada etapa (mudança de base
triangular). Para $H$ separável: de uma sequência densa, extraia uma
subfamília linearmente independente com subespaço gerado denso
(descarte cada vetor que esteja no subespaço gerado por seus
antecessores --- o subespaço gerado não muda) e ortonormalize: o
resultado é total.
\end{proof}

\begin{theorem}[O sistema trigonométrico; Parseval, enfim]
\label{thm:b3:hilbert:fourier}
Em $L^2(\intcc{-\pi}\pi)$ com $\langle f, g\rangle =
\frac1{2\pi}\int_{-\pi}^\pi \bar fg$, a família $e_n(t) =
\eu^{\iu nt}$, $n \in \Z$, é uma \omterm{def:b3:hilbert:onb}{base hilbertiana}.
Consequentemente, para \emph{toda} $f \in L^2$ --- em particular, para toda
$f$ $2\pi$-periódica \omterm{def:b3:topology:continuity}{contínua} por partes --- com $c_n(f) =
\frac1{2\pi}\int_{-\pi}^{\pi}f(t)\eu^{-\iu nt}\dd t$:
\[
f = \sum_{n\in\Z}c_n(f)\,\eu^{\iu nt} \ \ \text{em } L^2,
\qquad
\frac1{2\pi}\int_{-\pi}^{\pi}\abs f^2 =
\sum_{n\in\Z}\abs{c_n(f)}^2 .
\]
Isso demonstra, em plena generalidade, a identidade de \omterm{thm:b3:hilbert:parseval}{Parseval} que o
segundo ano admitiu.\index{identidade de Parseval (geral)}
\end{theorem}

\begin{proof}
A ortonormalidade é um cálculo direto (segundo ano). Totalidade: seja
$f \in L^2$ $\perp$ a todos os $e_n$, isto é, com todos os
coeficientes de Fourier nulos. As funções \omterm{def:b3:topology:continuity}{contínuas}
$2\pi$-periódicas são densas em $L^2(\intcc{-\pi}\pi)$: com efeito, $\mathcal
C_c(\intoo{-\pi}\pi)$ é denso
(\cref{thm:b3:lp:density}(2)) e tais funções se estendem periódica e
\omterm{def:b3:topology:continuity}{continuamente}. Os polinômios trigonométricos são densos, para $\norm\cdot_\infty$,
entre as funções \omterm{def:b3:topology:continuity}{contínuas} periódicas (Stone--Weierstrass,
\cref{cor:b3:complete:weierstrass}(c)), e $\norm\cdot_2 \leq \norm\cdot_\infty$: os polinômios trigonométricos são densos em
$L^2$. Mas $f \perp$ a todo polinômio trigonométrico, logo $f \perp$ a um
subespaço denso: $f
\in (\text{denso})^\perp = \{0\}$ (\cref{thm:b3:hilbert:decomposition}). O critério (iv)
do~\cref{thm:b3:hilbert:parseval} conclui; e (ii) e (iii) se desdobram na fórmula
exibida (reindexando o conjunto enumerável $\Z$; a série de duas
pontas converge incondicionalmente --- as somas parciais sobre qualquer
família exaustiva convergem, pelo argumento da cauda $\ell^2$).
\end{proof}

\begin{theorem}[Lax--Milgram]\label{thm:b3:hilbert:laxmilgram}
Sejam $H$ um \omterm{def:b3:hilbert:inner}{espaço de Hilbert} real e $a \colon H\times H \to
\R$ bilinear,
\emph{\omterm{def:b3:topology:continuity}{contínua}} ($\abs{a(u,v)} \leq M\norm
u\norm v$) e \emph{coerciva} ($a(u, u) \geq \alpha\norm u^2$, $\alpha > 0$). Então, para
todo $\varphi \in H'$, existe um único $u \in H$ com\index{teorema de Lax--Milgram}
\[
a(u, v) = \varphi(v) \qquad \text{para todo } v \in H .
\]
\end{theorem}

\begin{proof}
Para $u$ fixado, $v \mapsto a(u, v)$ é linear \omterm{def:b3:topology:continuity}{contínua}: Riesz dá um único $Au \in H$
com $a(u,v) = \langle Au,
v\rangle$; $A$ é linear com $\norm{Au} \leq M\norm u$ (unicidade dos representantes e,
depois, a cota). Coercividade: $\alpha\norm u^2 \leq a(u,u) = \langle Au, u\rangle \leq
\norm{Au}\norm u$, de modo que $\norm{Au} \geq \alpha\norm u$: $A$ é
injetor com imagem fechada (uma sequência imagem $Au_n$ de Cauchy
força $u_n$ de Cauchy). A imagem é densa: $w \perp
\operatorname{im}A$ dá $0 = \langle Aw, w\rangle \geq
\alpha\norm w^2$. Fechada e
densa: $A$ é bijetor. Dado $\varphi$, seja $f$ seu representante
(Riesz) e $u = A^{-1}f$: $a(u, v) = \langle f, v\rangle = \varphi(v)$, de maneira única ($a(u - u', \cdot) = 0$ e coercividade).
\end{proof}

\begin{remark}\label{rem:b3:hilbert:variational}
Quando $a$ é simétrica, a solução de Lax--Milgram é o único
minimizador da \emph{energia} $J(v) = \frac12a(v,v) -
\varphi(v)$ (\cref{exo:b3:hilbert:9}): a existência de
soluções de problemas variacionais de uma só vez. Aplicado a espaços de
funções adequados (os espaços de Sobolev de um curso posterior), isso
resolve problemas de valores de contorno para equações diferenciais ---
a porta de entrada moderna às equações diferenciais parciais.
\end{remark}

\section{Exercícios}

\begin{exercise}[$\star$]\label{exo:b3:hilbert:1}
(a) Demonstre as identidades de polarização (caso real: $4\langle x,
y\rangle = \norm{x+y}^2 - \norm{x-y}^2$; caso
complexo: a versão com quatro termos).
(b) Mostre que $\norm\cdot_1$ em $L^1(\intcc01)$ e $\norm\cdot_\infty$ em $\mathcal C(\intcc01)$ violam a lei do
paralelogramo: essas normas não provêm de \omterm{def:b3:hilbert:inner}{produto interno} algum.
\end{exercise}

\begin{exercise}[$\star$]\label{exo:b3:hilbert:2}
Em $H = L^2(\intcc01)$ (real):
(a) calcule a projeção de $f$ sobre o subespaço das funções
constantes e interprete;
(b) calcule a projeção sobre $\{g : g = 0 \text{ q.t.p.\ em }
\intcc0{1/2}\}$;
(c) calcule $d\bigl(x \mapsto x,\ \operatorname{Vect}(\mathbf
1)\bigr)$.
\end{exercise}

\begin{exercise}[$\star\star$]\label{exo:b3:hilbert:3}
(a) Mostre que, para um subespaço $F$: $F$ denso $\iff$
$F^\perp =
\{0\}$, e dê um exemplo em $\ell^2$ de um subespaço denso
\emph{próprio} (de modo que $F^\perp = 0$ sem que $F = H$: o teorema da
decomposição precisa mesmo de $F$ fechado).
(b) Mostre que, se $x_n \to x$ e $y_n \to y$ em norma, então $\langle x_n, y_n\rangle \to \langle x, y\rangle$, e localize dois
pontos em que o capítulo usou essa \omterm{def:b3:topology:continuity}{continuidade}.
\end{exercise}

\begin{exercise}[$\star\star$]\label{exo:b3:hilbert:4}
Aplique Gram--Schmidt a $1, x, x^2$ em $L^2(\intcc{-1}1)$ (\omterm{def:b3:measure:lebesgueouter}{medida de Lebesgue}): obtenha os
três primeiros \emph{polinômios de Legendre} normalizados e verifique
que coincidem com $\sqrt{n + \frac12}\,P_n$ para os polinômios de Rodrigues $P_n$
do~\cref{pb:b3:hilbert:1}.
\end{exercise}

\begin{exercise}[$\star\star$]\label{exo:b3:hilbert:5}
Aplique \omterm{thm:b3:hilbert:parseval}{Parseval} (\cref{thm:b3:hilbert:fourier}) a $f(t) = t$ e a $f(t) = t^2$ em $\intcc{-\pi}\pi$ --- agora
legitimamente para essas funções (\omterm{def:b3:topology:continuity}{contínuas}, mas cuja identidade
requeria antes cuidados de tipo $\mathcal C^1$ por partes na descontinuidade do
emenda): recupere
\[
\sum_{n\geq1}\frac1{n^2} = \frac{\pi^2}6,
\qquad
\sum_{n\geq1}\frac1{n^4} = \frac{\pi^4}{90} .
\]
\end{exercise}

\begin{exercise}[$\star\star$]\label{exo:b3:hilbert:6}
(a) Encontre $a \in L^2(\intcc01)$ com $\int_0^{1/2}f =
\langle a, f\rangle$ para toda $f$; calcule $\norm\varphi$ para esse
funcional.
(b) Mostre que a avaliação $f \mapsto f(\frac12)$, definida no subespaço $\mathcal C(\intcc01) \subseteq
L^2(\intcc01)$,
\emph{não} é \omterm{def:b3:topology:continuity}{contínua} para $\norm\cdot_2$: não existe representante de Riesz (a
avaliação não é uma noção $L^2$).
\end{exercise}

\begin{exercise}[$\star\star\star$]\label{exo:b3:hilbert:7}
Sejam $H$ separável com \omterm{def:b3:hilbert:onb}{base hilbertiana} $(e_n)$ e $(x_k)$
uma sequência limitada.
(a) Mostre que alguma subsequência converge \emph{fracamente}: existe
$x$ com $\langle y, x_{k_j}\rangle \to \langle y,
x\rangle$ para todo $y \in H$. \emph{(Extração diagonal nos
coeficientes $\langle e_n, x_k\rangle$; monte $x$ via Bessel e a limitação uniforme das
normas.)}
(b) Mostre que $e_n \rightharpoonup 0$, mas $\norm{e_n} = 1$: os limites fracos podem perder norma.
Mostre que $\norm x \leq
\liminf\norm{x_{k_j}}$ em (a).
\end{exercise}

\begin{exercise}[$\star\star$]\label{exo:b3:hilbert:8}
(Adjuntos) Para $T \in \mathcal L(H)$, mostre que existe um único $T^* \in \mathcal L(H)$ com $\langle Tx, y\rangle = \langle
x, T^*y\rangle$
(Riesz), e que $\vertiii{T^*} = \vertiii T$. Calcule o adjunto do deslocamento $S$ em
$\ell^2$ e demonstre $\ker T^* = (\operatorname{im}T)^\perp$ --- deduza $\overline{\operatorname{im}T} = (\ker T^*)^\perp$.
\end{exercise}

\begin{exercise}[$\star\star$]\label{exo:b3:hilbert:9}
Seja $a$ como em Lax--Milgram e, além disso, \emph{simétrica}.
Mostre que $u$ resolve $a(u, \cdot) = \varphi$ se, e somente se, $u$ minimiza
$J(v) = \frac12a(v, v) - \varphi(v)$, e que o mínimo é atingido em exatamente um ponto. \emph{(Complete
o quadrado: $J(u + w) - J(u) = \frac12a(w,w) \geq
\frac\alpha2\norm w^2$.)} Aplicação: redemonstre o teorema da projeção sobre
subespaços fechados a partir de Lax--Milgram.
\end{exercise}

\begin{exercise}[$\star\star\star$]\label{exo:b3:hilbert:10}
(O sistema de Haar) Em $\intcc01$, seja $h_{0} = \mathbf 1$ e, para $n = 2^j + k$ ($j \geq 0$, $0 \leq k < 2^j$):
\[
h_n = 2^{j/2}\Bigl(\mathbf 1_{[k2^{-j},\,(k +
\frac12)2^{-j})} - \mathbf 1_{[(k+\frac12)2^{-j},\,(k+1)2^{-j})}
\Bigr).
\]
Mostre que $(h_n)_{n\geq0}$ é ortonormal em $L^2(\intcc01)$ e total. \emph{(Ortogonalidade:
suportes disjuntos ou encaixados; totalidade: os subespaços gerados
finitos contêm todas as funções escada diádicas, que são densas --- via
o~\cref{thm:b3:lp:density}(1) e aproximação diádica de intervalos.)} O sistema de
Haar é o ancestral das ondaletas.
\end{exercise}

\begin{exercise}[$\star\star$]\label{exo:b3:hilbert:11}
(Projeções ortogonais, caracterizadas) Sejam $H$ um \omterm{def:b3:hilbert:inner}{espaço de
Hilbert} e $P \in \mathcal L(H)$ com $P^2 = P$, $P \neq 0$. Mostre a equivalência de:
(i) $P$ é a projeção ortogonal sobre $\operatorname{im}P$;
(ii) $P = P^*$ (\cref{exo:b3:hilbert:8});
(iii) $\vertiii P = 1$.
\emph{(Para (iii) $\Rightarrow$ (i): se algum $x \in
(\ker P)^\perp$ tivesse $Px \neq x$, considere
$x + t(Px - x)$ --- ou, diretamente: para $u \in \operatorname{im}P$ e $v \in
\ker P$, expanda $\norm{P(u + tv)}^2 \leq \norm{u + tv}^2$ para todo
$t \in \R$ e conclua $\langle u, v\rangle = 0$.)}
Exiba uma projeção não ortogonal em $\R^2$ e calcule sua norma.
\end{exercise}

\begin{exercise}[$\star\star\star$]\label{exo:b3:hilbert:12}
(Teorema ergódico de von Neumann) Sejam $U \in \mathcal L(H)$ \emph{unitário} ($U^*U = UU^* = I$),
$F = \ker(U - I)$ o espaço dos pontos fixos, $P$ a projeção ortogonal sobre
$F$ e $A_n = \frac1n\sum_{k=0}^{n-1}U^k$.
(a) Mostre que $\ker(U - I) = \ker(U^* - I)$ \emph{(a partir de $\norm{Ux - x}^2 = 2\norm x^2 - 2\operatorname{Re}\langle
Ux, x\rangle$ e da unitariedade)} e
deduza $\overline{\operatorname{im}(U - I)} = F^\perp$.
(b) Mostre que $A_nx \to x$ para $x \in F$, e $A_nx \to 0$ para $x \in \operatorname{im}(U - I)$
\emph{(telescopamento)}, e depois para $x \in \overline{\operatorname{im}(U - I)}$ (cota uniforme $\vertiii{A_n} \leq 1$).
(c) Conclua: $A_nx \to Px$ para \emph{todo} $x \in H$ --- as médias temporais
convergem à projeção sobre os invariantes.
(d) Detalhe o caso $H = L^2(\R/\Z)$ e $Uf = f(\cdot +
\alpha)$ com $\alpha$ irracional: identifique
$F$ (use séries de Fourier, \cref{thm:b3:hilbert:fourier}) e deduza que $\frac1n\sum_{k<n}f(x + k\alpha) \to \int_0^1f$ em
$L^2$: a equidistribuição $L^2$ das rotações irracionais.
\end{exercise}

\section{Problema: polinômios ortogonais}

\begin{problem}[{Problema de fim de semana --- Legendre, Hermite e a
quadratura de Gauss}]\label{pb:b3:hilbert:1}
Sejam $I \subseteq \R$ um intervalo e $w > 0$ um \emph{peso} \omterm{def:b3:topology:continuity}{contínuo} no
\omterm{def:b3:topology:interior}{interior} de $I$ tal que $\int_I
\abs t^nw(t)\dd t < \infty$ para todo $n$; trabalhamos em
$H = L^2(I,
w\,\dd\lambda)$ com $\langle f, g\rangle = \int_I \bar
fg\,w$. Gram--Schmidt aplicado a $1, t, t^2, \dots$ produz os
\emph{\omterm{pb:b3:hilbert:1}{polinômios ortogonais}} $(p_n)$ para $w$ (com normalização
mônica: $p_n = t^n + \cdots$).\index{polinômios ortogonais}

\medskip
\noindent\textbf{Parte I --- Teoria geral.}
\begin{enumerate}
  \item Mostre que $p_n$ é ortogonal a todo polinômio de grau
        $< n$ e que $(p_0, \dots, p_n)$ é uma base de $\R_n[t]$.
  \item (Recorrência de três termos) Mostre que existem reais $a_n,
        b_n$ com
        \[
        p_{n+1}(t) = (t - a_n)\,p_n(t) - b_n\,p_{n-1}(t),
        \qquad b_n = \frac{\norm{p_n}^2}{\norm{p_{n-1}}^2} >
        0 .
        \]
        \emph{(Expanda $t\,p_n$ na base $(p_k)_{k \leq
        n+1}$ e anule coeficientes
        por ortogonalidade, usando $\langle tp_n, p_k\rangle = \langle p_n,
        tp_k\rangle$.)}
  \item (Raízes) Mostre que $p_n$ tem $n$ raízes \emph{distintas},
        todas \omterm{def:b3:topology:interior}{interiores} a $I$. \emph{(Sejam $t_1 < \dots < t_m$ as mudanças de
        sinal \omterm{def:b3:topology:interior}{interiores} de $p_n$; se $m < n$, teste $p_n$
        contra $\prod_{i\leq m}(t - t_i)$ e contradiga a ortogonalidade.)}
\end{enumerate}

\medskip
\noindent\textbf{Parte II --- Legendre ($I = \intcc{-1}1$, $w =
1$).} Defina
$P_n(t) = \frac{1}{2^nn!}\,\frac{\dd^n}{\dd
t^n}\bigl[(t^2 - 1)^n\bigr]$ (Rodrigues).
\begin{enumerate}[resume]
  \item Mostre que $\deg P_n = n$ com coeficiente líder $\frac{(2n)!}{2^n(n!)^2}$ e, integrando por
        partes $n$ vezes, que $\langle P_n, Q\rangle = 0$ para todo polinômio $Q$ de
        grau $< n$: os $P_n$ são (a menos de normalização) os
        \omterm{pb:b3:hilbert:1}{polinômios ortogonais} para $w = 1$.
  \item Calcule $\norm{P_n}_2^2 = \frac{2}{2n+1}$ \emph{(integre por partes $n$ vezes contra si
        mesmo e reduza a uma integral de Beta/Wallis,
        \cref{exo:b3:product:8})}.
  \item Mostre que os polinômios de Legendre normalizados formam uma
        \omterm{def:b3:hilbert:onb}{base hilbertiana} de $L^2(\intcc{-1}1)$ \emph{(Weierstrass,
        \cref{cor:b3:complete:weierstrass}, mais a \omterm{ex:b3:lebesgue:gamma}{densidade} de $\mathcal C$ em $L^2$)}, e
        expanda $f(t) = \abs t$ até o grau $2$: calcule a melhor aproximação
        quadrática em $L^2$ de $\abs t$.
\end{enumerate}

\medskip
\noindent\textbf{Parte III --- Hermite ($I = \R$, $w(t) =
\eu^{-t^2}$).} Defina
$H_n(t) =
(-1)^n\eu^{t^2}\frac{\dd^n}{\dd t^n}\eu^{-t^2}$.
\begin{enumerate}[resume]
  \item Mostre que $H_n$ é um polinômio de grau $n$ com
        coeficiente líder $2^n$, que $H_{n+1} = 2tH_n -
        H_n'$ e que $\langle H_m, H_n\rangle_w =
        \delta_{mn}\,2^nn!\sqrt\pi$ \emph{(partes,
        de novo)}.
  \item Mostre que a família de Hermite é total em $L^2(\R,
        \eu^{-t^2}\dd t)$, admitindo um
        resultado do~\cref{ch:b3:fouriertransform}: se $g \in L^1(\R)$ satisfaz $\int g(t)\eu^{-\iu\xi t}\dd t = 0$ para todo
        $\xi$, então $g = 0$ q.t.p. \emph{(Para $f \perp$ a todos os
        $H_n$, isto é, $\perp$ a todos os polinômios: mostre que
        $z \mapsto \int
        f(t)\eu^{-t^2}\eu^{-\iu zt}\dd t$ está bem definida, expanda a exponencial em série,
        justifique a troca por dominação e conclua que a transformada
        de Fourier de $f\eu^{-t^2}$ se anula.)}
\end{enumerate}

\medskip
\noindent\textbf{Parte IV --- Quadratura de Gauss.} Fixe $n$, sejam
$t_1 < \dots < t_n$ as raízes de $p_n$ (Parte I) e defina os pesos $w_i = \int_I \ell_i(t)\,w(t)\dd t$, em que
$\ell_i$ são os polinômios da base de interpolação de Lagrange nos
$t_i$.
\begin{enumerate}[resume]
  \item Mostre que a regra de quadratura $Q(f) = \sum_iw_if(t_i)$ é exata em todos os
        polinômios de grau $\leq n - 1$ (interpolação) e, de fato --- eis o
        milagre ---, em todos os polinômios de grau $\leq 2n - 1$: escreva
        $P =
        qp_n + r$ e use a ortogonalidade sobre o quociente $q$.
  \item Mostre que os pesos são positivos \emph{(aplique a regra a
        $\ell_i^2$, de grau $2n -
        2$)} e deduza do teorema de Pólya
        (\cref{exo:b3:banach:9}) que a quadratura de Gauss converge: $Q_n(f) \to \int_I fw$ para
        toda $f$ \omterm{def:b3:topology:continuity}{contínua} em um $I$ \omterm{def:b3:topology:compact}{compacto}.
  \item Para $n = 2$, $I = \intcc{-1}1$, $w = 1$: calcule os nós $\pm\frac1{\sqrt3}$ e os
        pesos $1, 1$, e verifique a exatidão em $1, t, t^2, t^3$ à mão. Compare
        com a regra do trapézio nos mesmos dois pontos de avaliação.
\end{enumerate}

\medskip
\noindent\textbf{Parte V --- Chebyshev: os polinômios que melhor
oscilam.} Agora $I = \intcc{-1}1$ e $w(t) =
\frac1{\sqrt{1 - t^2}}$.
\begin{enumerate}[resume]
  \item Mostre que $T_n(\cos\theta) = \cos n\theta$ define um polinômio $T_n$ de grau $n$
        (estabeleça $T_{n+1} =
        2t\,T_n - T_{n-1}$ a partir de uma identidade trigonométrica),
        com coeficiente líder $2^{n-1}$ para $n \geq 1$; e que a substituição
        $t = \cos\theta$ dá
        \[
        \langle T_m, T_n\rangle_w =
        \int_0^\pi\cos m\theta\,\cos n\theta\,\dd\theta
        = 0 \ (m \neq n), \qquad
        \norm{T_0}_w^2 = \pi,\ \ \norm{T_n}_w^2 = \frac\pi2 :
        \]
        os $T_n$ são os \omterm{pb:b3:hilbert:1}{polinômios ortogonais} para esse peso, e as
        expansões de Chebyshev \emph{são} séries de Fourier em cosseno
        disfarçadas.
  \item Localize explicitamente as $n$ raízes $t_k =
        \cos\frac{(2k-1)\pi}{2n}$ e os $n + 1$
        extremos $s_j = \cos\frac{j\pi}n$ de $T_n$ em $\intcc{-1}1$, em que $T_n(s_j) = (-1)^j$: o gráfico
        \emph{equioscila} entre $\pm1$.
  \item (Minimax) Mostre que, entre todos os polinômios \emph{mônicos}
        de grau $n$, o polinômio $2^{1-n}T_n$ tem a menor norma do supremo
        em $\intcc{-1}1$, a saber $2^{1-n}$ --- e que ele é o único minimizador.
        \emph{(Se um $P$ mônico tivesse $\sup\abs P < 2^{1-n}$, a diferença $2^{1-n}T_n -
        P$, de
        grau $\leq n-1$, alternaria de sinal nos $n+1$ pontos de
        equioscilação.)}
  \item Aplicação à interpolação: para nós $t_1 < \dots
        < t_n$ em $\intcc{-1}1$, o erro da
        interpolação de Lagrange de uma função $\mathcal C^n$ envolve $\omega(t) = \prod_i(t - t_i)$.
        Mostre que escolher as raízes de Chebyshev como nós minimiza
        $\sup_{\intcc{-1}1}\abs\omega$, e dê a cota resultante $\norm{f -
        L_nf}_\infty \leq \frac{\norm{f^{(n)}}_\infty}
        {2^{n-1}\,n!}$ --- compare com nós
        igualmente espaçados (enuncie o fenômeno de Runge como o conto
        de advertência).
  \item Verifique $\abs{T_n'(\pm1)} = n^2$ \emph{(derive $T_n(\cos\theta) = \cos n\theta$ e tome os limites $\theta \to 0, \pi$)}: um
        polinômio limitado por $1$ em $\intcc{-1}1$ pode ter derivada tão
        grande quanto $n^2$ na borda (a desigualdade de Markov diz
        que não pode ser maior --- apenas o enunciado). Em que parte do
        intervalo a cota da derivada é apenas $O(n)$?
  \item (Quadratura de Chebyshev--Gauss) Mostre que a regra de Gauss
        para o peso $w$ nas $n$ raízes de Chebyshev tem pesos
        \emph{iguais} $w_i = \frac\pi n$ \emph{(exatidão em $T_0, \dots, T_{n-1}$ mais as somas
        trigonométricas $\sum_{k=1}^n\cos\bigl(j\tfrac{(2k-1)\pi}{2n}\bigr) =
        0$ para $1 \leq j \leq n - 1$)}: a mais uniforme de todas as
        quadraturas. Escreva-a por extenso para $n = 3$.
\end{enumerate}

\medskip
\noindent\textbf{Parte VI --- Christoffel--Darboux, entrelaçamento e a
matriz de Jacobi.} De volta a um peso geral; $h_k = \norm{p_k}^2$ (com $p_k$
mônicos), $b_k = h_k/h_{k-1}$.
\begin{enumerate}[resume]
  \item (Norma mínima) Mostre que, entre todos os polinômios
        \emph{mônicos} de grau $n$, o ortogonal $p_n$ é o único
        de norma $L^2(w)$ mínima --- identifique a minimização como
        uma projeção ortogonal sobre $\R_{n-1}[t]$ (\cref{thm:b3:hilbert:projection} ou a projeção
        em dimensão finita do segundo ano). A propriedade minimax da
        questão 14 é o mesmo enunciado com $L^\infty$ no lugar de $L^2$:
        mesmo herói, duas normas.
  \item (Christoffel--Darboux) Demonstre, por indução sobre $n$
        usando a recorrência de três termos, a identidade
        \[
        \sum_{k=0}^{n}\frac{p_k(x)\,p_k(y)}{h_k}
        = \frac{p_{n+1}(x)\,p_n(y) -
        p_n(x)\,p_{n+1}(y)}{h_n\,(x - y)}
        \qquad (x \neq y),
        \]
        e sua forma confluente ($y \to x$): $\sum_{k\leq n}\frac{p_k(x)^2}{h_k} =
        \frac{p_{n+1}'(x)p_n(x) - p_n'(x)p_{n+1}(x)}{h_n}$.
  \item Deduza que $p_n$ e $p_{n+1}$ não têm raiz comum e que, em toda
        raiz $x_0$ de $p_{n+1}$: $p_n(x_0)\,p_{n+1}'(x_0) > 0$. Conclua o \emph{entrelaçamento}
        das raízes: entre duas raízes consecutivas de $p_{n+1}$ há
        exatamente uma raiz de $p_n$.
  \item (Matriz de Jacobi) Seja $J_n$ a matriz simétrica
        tridiagonal $n\times n$ de diagonal $a_0,
        \dots, a_{n-1}$ e entradas fora da diagonal
        $\sqrt{b_1},
        \dots, \sqrt{b_{n-1}}$. Mostre por indução que $\det(tI_n - J_n) = p_n(t)$, de modo que as raízes de
        $p_n$ são os autovalores de uma matriz real simétrica ---
        redemonstrando em uma linha que elas são reais e, com o
        entrelaçamento acima, amarrando os \omterm{pb:b3:hilbert:1}{polinômios ortogonais} ao
        mundo espectral do~\cref{ch:b3:spectral}.
  \item (Síntese) Monte o dicionário das três famílias clássicas
        (Legendre, Hermite, Chebyshev): intervalo, peso, fórmula de
        definição, recorrência de três termos, norma e o habitat
        natural de cada uma (quadratura e aproximação em \omterm{def:b3:topology:compact}{compactos};
        análise gaussiana; métodos minimax e de Fourier em cosseno).
        Uma frase sobre o que a teoria geral (Partes I e VI) deu e que
        nenhum cálculo individual poderia dar.
\end{enumerate}

\medskip
\noindent\textbf{Parte VII --- O termo de erro e o núcleo por trás dos
pesos.} Aqui $I$ é \omterm{def:b3:topology:compact}{compacto} e $f \in \mathcal
C^{2n}(I)$.
\begin{enumerate}[resume]
  \item (Fórmula do erro de Gauss) Seja $Hf$ o interpolante de
        Hermite de grau $\leq 2n - 1$ que coincide com $f$ e $f'$ nos nós
        $t_1, \dots, t_n$ (demonstre sua existência e o erro pontual
        \[
        f(t) - Hf(t) =
        \frac{f^{(2n)}(\xi_t)}{(2n)!}\;p_n(t)^2
        \]
        pelo argumento habitual da função auxiliar). Deduza,
        integrando essa identidade contra $w$ e confinando entre os
        extremos de $f^{(2n)}$, que
        \[
        \int_I f\,w - Q_n(f)
        = \frac{f^{(2n)}(\xi)}{(2n)!}\,h_n
        \qquad\text{para algum } \xi \in I,
        \]
        com $h_n = \norm{p_n}^2$ como na Parte VI: a quadratura de Gauss erra por uma
        derivada de ordem $2n$, ponderada pela norma ao quadrado do
        polinômio ortogonal mônico.
  \item (Os pesos são valores de Christoffel) Usando o núcleo
        reprodutor $K_n(x, y) = \sum_{k=0}^{n-1}
        \frac{p_k(x)p_k(y)}{h_k}$ de $\R_{n-1}[t]$ e a exatidão de $Q_n$ até o grau
        $2n - 2$, demonstre
        \[
        w_i \;=\;
        \Bigl(\,\sum_{k=0}^{n-1}
        \frac{p_k(t_i)^2}{h_k}\Bigr)^{\!-1} :
        \]
        cada peso é o valor, em seu nó, da \emph{função de
        Christoffel} --- a positividade dos pesos (questão 10) de
        novo, agora com fórmula exata. Verifique que ela recupera $w_1 = w_2 = 1$
        para $n =
        2$, $I = \intcc{-1}1$, $w = 1$.
  \item (Tudo se confere em uma integral) Para o peso de Chebyshev e
        $n = 3$ nós, calcule os dois lados de
        \[
        \int_{-1}^{1}\frac{t^6}{\sqrt{1 - t^2}}\,\dd t
        = \frac{5\pi}{16},
        \qquad
        Q_3(t^6) = \frac{9\pi}{32},
        \]
        de modo que o erro de quadratura é exatamente $\frac{\pi}{32}$; verifique
        então que a fórmula de erro da questão 23 prevê exatamente esse
        valor (aqui $f^{(6)} = 6!$ é constante, e $h_3 = \norm{2^{-2}T_3}_w^2 =
        \frac\pi{32}$): teoria e cálculo
        concordam até o último algarismo.
\end{enumerate}
\end{problem}
